\section{Discussion}
\label{sec:discussion}

Several protocol choices changed the reported error while the trained weights
remained fixed. In the dataset-transfer experiment, the annotation convention
changed with the dataset. The synthetic-occlusion result depended partly on the
geometry and appearance of the occluder. A count produced at a fixed threshold
also reflected the calibration of the model's confidence scores. The
photometric sweep measured the response to simulated camera settings. The
deployment comparison included changes in precision, NMS, padding, silicon, and
operating-system version. The metrics remain useful when the protocol specifies
these factors and reports the relevant controls and complementary measurements.

\subsection{Capacity has the smallest measured effect}

The measured effect sizes differ substantially. Moving from the smallest to the
largest model on this split changes clean recall by 0.058 and mAP50--95 by
0.068. The improvement ends at the medium model, two sizes below the largest
model tested.
Over-exposing the same frames by three stops reduces recall by 0.15 to 0.21,
while covering 40\% of each fruit reduces it by more than 0.3. Changing the
annotation convention changes the assigned stage for a large fraction of
correctly detected half-ripened fruit. Selecting a different confidence
threshold changes cross-domain counting error by a factor of two without
changing mAP.

In this study, threshold calibration and exposure control changed recall or
counting error more than model scaling. Model capacity is nevertheless the
variable most often optimised. This comparison concerns measured effect sizes;
the study did not measure the implementation cost of the interventions.

\subsection{What the fully-ripe count says about capacity}

Detection improves across part of the model-size series, but fully-ripe counting
error does not. The smallest model has the best value
(Section~\ref{sec:ladder}). A stage count requires the detector to find the
fruit and then assign a maturity stage from the colour within the detected box.
The stage assignment is affected by the annotation boundary in
Section~\ref{sec:colour} and the illumination in Section~\ref{sec:lighting}.
Increasing the network's depth and width does not address either factor.
Figure~\ref{fig:detection-examples} shows an example: on the same frame the
largest model labels a clearly red fruit half-ripened while the smallest model
agrees with the annotation. Fully-ripe MAE should therefore be reported
alongside mAP.

\subsection{Latency depends on the deployment device}

The laptop-to-smartphone latency ratio falls from 3.85$\times$ for the smallest
model to 1.21$\times$ for the largest. A development-machine benchmark therefore
underestimates smartphone latency most strongly for the small models considered
for mobile deployment. The measurement also changes over time. An
operating-system update changed the smallest model's accelerator latency by
67\% while leaving a larger model unchanged. Reproducing a latency result
requires the device, compute-unit setting, which specifies the permitted
processor combination, OS version, and run duration.

\subsection{Augmentation ablations need an out-of-distribution test}

A clean-only ablation would favour removing the colour jitter. Removing it has
no measured penalty on clean frames and gives the best clean mAP among the four
recipes. Under a perturbation generated by a different process, however, the
model trained without colour jitter loses up to 0.095 recall. Claims about
augmentation concern behaviour outside the training distribution and should be
evaluated under such conditions. The test intervention should also differ from
the augmentation process used in training.

The mixed-data result requires a related qualification. The two source domains
are visually distinguishable, so the network can learn a different maturity
boundary for each source. Random label flips provide no corresponding visual
cue. This distinction motivates a broader study comparing label noise whose
source can be inferred from the image with noise whose source cannot be
inferred.

The occlusion study shows the value of testing several levels of severity,
consistent with prior controlled benchmarks \cite{gilroy2023benchmark}.
It also confirms the need for multiple occluders and natural controls. The gap
between the background and black patches shows the effect of occluder
appearance. The smaller loss under untouched fruit overlap also agrees with
evidence that robustness to artificial occlusion need not transfer to realistic
occlusion \cite{kassaw2026occlusion}.

\subsection{Comparison with previous work}

Previous tomato maturity detectors report mAP50 values between about 0.86 and
0.97 \cite{su2022seyolov3, li2023yolov5tomato, li2023mhsayolov8,
camacho2023tomato}, each on its own dataset, with its own class definitions and
its own train/test split. These differences in evaluation conditions prevent a
direct comparison with one another or with the values reported here. In
Table~\ref{tab:cross}, the same YOLO11n detector, trained once, scores 0.698
mAP50--95 on the dataset it was trained on and 0.389 on the other public tomato
dataset, which uses different annotators and a different colour boundary for
the ripe stages. A difference of a few points between two published detectors
is much smaller than the difference produced by changing the test set alone.

The datasets also have different policies for unlabelled fruit. tomatOD leaves
fruit that are more than half hidden unannotated \cite{tsironis2020tomatod},
so a detector evaluated on it is never asked to find heavily occluded fruit,
whereas Section~\ref{sec:occlusion-results} shows that recall falls steeply
above 50\% hidden. Studies that target occluded fruit directly, such as Mu
\emph{et al.}\ for green tomatoes \cite{mu2020occluded}, therefore measure a
different task from studies that report accuracy on tomatOD, even when both
report mAP.

On lighting, the results agree with Haggag \emph{et al.}, who found
illumination the hardest of the shifts they tested for a greenhouse tomato
detector \cite{haggag2024generalization}. Here, three stops of over-exposure
reduced recall by more than three times as much as the whole difference between
the smallest and the largest model. Unlike that study, this work also tested
whether the colour augmentation already in the training pipeline explains the
measured robustness, and repeated the comparison over three seeds.

On deployment, earlier tomato studies on smartphones and embedded boards
report the speed of one model on one platform \cite{zeng2023mobile, jia2025tseyolo,
liu2025gaeyolo, kafi2025drone}. The model-size series in
Section~\ref{sec:ladder} shows
that the ratio between laptop and smartphone latency changes from 3.85 to 1.21
across model sizes, and that an operating-system update changed the latency of
one model by 67\% while leaving another unchanged. A single speed figure from
another device, model size or operating-system version is therefore a weak
guide to the speed of a different configuration. A deployment study needs to
report all four: the device, model size, compute-unit setting and
operating-system version.
